\clearpage
\chapter*{Fe de erratas de la exposición n.\textsuperscript{o} 232}
\addcontentsline{toc}{chapter}{Fe de erratas de la exposición n.\textsuperscript{o} 232}
\src{303}

\noindent\textbf{7. TDTE V: Esquemas de Picard. Teoremas de existencia.}

\begin{center}
[Seminario Bourbaki, vol. 14, 1961/62, n.\textsuperscript{o} 232, 19 pp.]
\end{center}

\noindent\textbf{Página 7, observación 3.3.} --- La cuestión planteada aquí fue
resuelta afirmativamente por MUMFORD (véanse los Comentarios a la exposición
siguiente).

\noindent\textbf{Página 13, observación 5.2.} --- Como señalo al comienzo de la
exposición siguiente, la conjetura de existencia formulada aquí es falsa; no
obstante, MUMFORD consigue demostrar por sus métodos un teorema ligeramente más
débil.

\noindent\textbf{Página 14, línea 7.} --- Después de «completo», añádase «con
cuerpo residual algebraicamente cerrado». El contraejemplo de MUMFORD recién
mencionado muestra, además, que esta restricción es indispensable.

\noindent\textbf{Página 17, observación 6.6.} --- Como señalamos al comienzo de
la exposición siguiente, la cuestión planteada aquí acaba de ser resuelta
afirmativamente por MURRE.

\noindent\textbf{Página 18, observación 6.7.} --- La cuestión planteada aquí
tiene respuesta afirmativa (cf. el último párrafo de los Comentarios a la
exposición siguiente).

\noindent\textbf{Página 18, observación 6.8, cuarta línea.} --- En lugar de
«regular», léase «liso sobre $k$».
